% answers_short.tex —— 2006 年·普通高等学校招生考试·福建卷文科 · 简短答案速查
% 由 tools/make_short_answers.py 自 answers.tex 的 \gksolhead/\gkanswer 提取，
% 并与 2006/2006_数学真题_福建文_答案要点.md 逐题比对后生成，勿手改（改源 MD 后重跑）。

\begin{gknotice}
\textbf{使用说明}：本节只给最终答案便于快速核对，完整推导见解析版。答案为编写者独立推导，\textbf{非官方标准答案}。
\end{gknotice}

\gkpart{选择题}
\begin{center}\begin{tabular}{ccccccccccccc}
\toprule
题号 & 1 & 2 & 3 & 4 & 5 & 6 & 7 & 8 & 9 & 10 & 11 & 12 \\
\midrule
答案 & D & B & B & A & C & A & D & B & B & C & C & D \\
\bottomrule
\end{tabular}\end{center}

\gkpart{填空题}
\begin{itemize}[leftmargin=2.6em]
  \item[\textbf{13.}] $10$
  \item[\textbf{14.}] $\dfrac14$
  \item[\textbf{15.}] $4$
  \item[\textbf{16.}] $\dfrac32$
\end{itemize}

\gkpart{解答题}
\begin{itemize}[leftmargin=2.6em]
  \item[\textbf{17.}] (1) 最小正周期 $\pi$，单调递增区间 $\left[-\dfrac\pi3+k\pi,\dfrac\pi6+k\pi\right]$，$k\in\mathbb Z$；(2) 将 $y=\sin2x$ 的图象向左平移 $\dfrac\pi{12}$，再向上平移 $\dfrac32$
  \item[\textbf{18.}] (1) $\dfrac56$；(2) $\dfrac5{36}$；(3) $\dfrac5{16}$
  \item[\textbf{19.}] (1) 证明见解析；(2) $\arccos\dfrac1{2\sqrt2}$；(3) $\sqrt{\dfrac37}$
  \item[\textbf{20.}] (1) $\left(x+\dfrac12\right)^2+(y-\sqrt2)^2=\dfrac94$ 或 $\left(x+\dfrac12\right)^2+(y+\sqrt2)^2=\dfrac94$；(2) $x-2y+1=0$
  \item[\textbf{21.}] (1) $f(x)=2x^2-10x$；(2) $m=3$
  \item[\textbf{22.}] (1) 数列 $\{a_{n+1}-a_n\}$ 首项为 $2$、公比为 $2$；(2) $a_n=2^n-1$；(3) $\{b_n\}$ 为等差数列，可写为 $b_n=2+(n-1)(b_2-2)$
\end{itemize}
